\begin{examplebox}
	\textbf{\textcolor{CExample}{例 1（基础）}}
	\textbf{\textcolor{CExample}{题目：}}
	已知点$P(1,2)$和直线$l:2x+3y-1=0$，求$P$关于$l$的对称点$P'$的坐标。
	\textbf{\textcolor{CExample}{解题步骤：}}
	\begin{description}[style=nextline, leftmargin=2.8em, labelsep=0.5em, itemsep=0.45em]
		\item[\textbf{第一步（识别系数）：}]
			从直线方程得$A=2,B=3,C=-1$，点坐标$x_0=1,y_0=2$。
			\par\textbf{理由：} 公式需要对应系数。
		\item[\textbf{第二步（计算代数值）：}]
			计算$Ax_0+By_0+C=2\times1+3\times2-1=7$，分母$A^2+B^2=13$。
			\par\textbf{理由：} 这些值将代入公式。
		\item[\textbf{第三步（套用公式）：}]
			\[
				x'=1-\dfrac{2\times2\times7}{13}=1-\dfrac{28}{13}=-\dfrac{15}{13},\quad y'=2-\dfrac{2\times3\times7}{13}=2-\dfrac{42}{13}=-\dfrac{16}{13}.
			\]
			\par\textbf{理由：} 直接应用对称点公式。
	\end{description}
	\textbf{\textcolor{CExample}{关键结论：}}
	对称点坐标为$\left(-\dfrac{15}{13},-\dfrac{16}{13}\right)$。

	\medskip
	\textbf{\textcolor{CExample}{例 2（稍变形）}}
	\textbf{\textcolor{CExample}{题目：}}
	已知点$Q(0,3)$和直线$m:y=2x+1$，求$Q$关于$m$的对称点$Q'$的坐标。
	\textbf{\textcolor{CExample}{解题步骤：}}
	\begin{description}[style=nextline, leftmargin=2.8em, labelsep=0.5em, itemsep=0.45em]
		\item[\textbf{第一步（化为一般式）：}]
			直线$y=2x+1$化为$2x-y+1=0$，故$A=2,B=-1,C=1$，点坐标$x_0=0,y_0=3$。
			\par\textbf{理由：} 公式要求直线为一般式。
		\item[\textbf{第二步（计算代数值）：}]
			计算$Ax_0+By_0+C=2\times0+(-1)\times3+1=-2$，分母$A^2+B^2=5$。
			\par\textbf{理由：} 代入公式所需。
		\item[\textbf{第三步（套用公式）：}]
			\[
				x'=0-\dfrac{2\times2\times(-2)}{5}=0+\dfrac{8}{5}=\dfrac{8}{5},\quad y'=3-\dfrac{2\times(-1)\times(-2)}{5}=3-\dfrac{4}{5}=\dfrac{11}{5}.
			\]
			\par\textbf{理由：} 直接应用对称点公式，注意符号。
	\end{description}
	\textbf{\textcolor{CExample}{关键结论：}}
	对称点坐标为$\left(\dfrac{8}{5},\dfrac{11}{5}\right)$。

	\medskip
	\textbf{\textcolor{CExample}{例 3（综合应用）}}
	\textbf{\textcolor{CExample}{题目：}}
	光线从点$A(1,0)$射向直线$l:x-y+1=0$，经反射后经过点$B(3,2)$，求入射点（反射点）$R$的坐标。
	\textbf{\textcolor{CExample}{解题步骤：}}
	\begin{description}[style=nextline, leftmargin=2.8em, labelsep=0.5em, itemsep=0.45em]
		\item[\textbf{第一步（求对称点）：}]
			求$A$关于$l$的对称点$A'$。直线$x-y+1=0$中$A=1,B=-1,C=1$，点$A(1,0)$。计算$Ax_0+By_0+C=1-0+1=2$，$A^2+B^2=2$。由公式得
			\[
				x'=1-\dfrac{2\times1\times2}{2}=1-2=-1,\quad y'=0-\dfrac{2\times(-1)\times2}{2}=0+2=2.
			\]
			故$A'(-1,2)$。
			\par\textbf{理由：} 反射光线所在直线过$A'$和$B$。
		\item[\textbf{第二步（求反射线方程）：}]
			由$A'(-1,2)$和$B(3,2)$得直线$A'B$的方程$y=2$（斜率$0$）。
			\par\textbf{理由：} 两点纵坐标相等。
		\item[\textbf{第三步（求交点）：}]
			解方程组$
			\begin{cases}
				y=2,\\x-y+1=0
		\end{cases}
			$得$x=1,y=2$，即$R(1,2)$。
			\par\textbf{理由：} $R$为入射点，也是反射点。
	\end{description}
	\textbf{\textcolor{CExample}{关键结论：}}
	反射点坐标为$(1,2)$。
\end{examplebox}
